% source: Robin Hartshorne, "Algebraic Geometry", GTM 52 (1977)
% pdf_page: 0148
% doc_page: 131 (Chapter II, Section 6, "Divisors")
% retrieved: 2026-07-26
% transcribed_by: codex/gpt-5 (vision, rendered at 200dpi via pdftoppm)

% running head: 6 Divisors / p. 131

\DeclareMathOperator{\Spec}{Spec}
\DeclareMathOperator{\Div}{Div}
\DeclareMathOperator{\Cl}{Cl}

% [continuation from page-0147.tex]

mined by its valuation (Ex. 4.5). Now let $f \in K^*$ be any nonzero rational
function on $X$. Then $v_Y(f)$ is an integer. If it is positive, we say $f$
has a \emph{zero} along $Y$, of that order; if it is negative, we say $f$
has a \emph{pole} along $Y$, of order $-v_Y(f)$.

\begin{lemma}[6.1]
Let $X$ satisfy $(*)$, and let $f \in K^*$ be a nonzero function on $X$.
Then $v_Y(f) = 0$ for all except finitely many prime divisors $Y$.
\end{lemma}

\begin{proof}
Let $U = \Spec A$ be an open affine subset of $X$ on which $f$ is regular.
Then $Z = X - U$ is a proper closed subset of $X$. Since $X$ is noetherian,
$Z$ can contain at most finitely many prime divisors of $X$; all the others
must meet $U$. Thus it will be sufficient to show that there are only
finitely many prime divisors $Y$ of $U$ for which $v_Y(f) \ne 0$. Since $f$
is regular on $U$, we have $v_Y(f) \ge 0$ in any case. And $v_Y(f) > 0$ if
and only if $Y$ is contained in the closed subset of $U$ defined by the ideal
$Af$ in $A$. Since $f \ne 0$, this is a proper closed subset, hence contains
only finitely many closed irreducible subsets of codimension one of $U$.
\end{proof}

\begin{definition}
Let $X$ satisfy $(*)$ and let $f \in K^*$. We define the \emph{divisor of
$f$}, denoted $(f)$, by
\[
  (f) = \sum v_Y(f) \cdot Y,
\]
where the sum is taken over all prime divisors of $X$. By the lemma, this is a
finite sum, hence it is a divisor. Any divisor which is equal to the divisor
of a function is called a \emph{principal divisor}.
\end{definition}

Note that if $f,g \in K^*$, then $(f/g) = (f) - (g)$ because of the
properties of valuations. Therefore sending a function $f$ to its divisor
$(f)$ gives a homomorphism of the multiplicative group $K^*$ to the additive
group $\Div X$, and the image, which consists of the principal divisors, is a
subgroup of $\Div X$.

\begin{definition}
Let $X$ satisfy $(*)$. Two divisors $D$ and $D'$ are said to be
\emph{linearly equivalent}, written $D \sim D'$, if $D - D'$ is a principal
divisor. The group $\Div X$ of all divisors divided by the subgroup of
principal divisors is called the \emph{divisor class group} of $X$, and is
denoted by $\Cl X$.
\end{definition}

The divisor class group of a scheme is a very interesting invariant. In
general it is not easy to calculate. However, in the following propositions
and examples we will calculate a number of special cases to give some idea of
what it is like.

\begin{proposition}[6.2]
Let $A$ be a noetherian domain. Then $A$ is a unique factorization domain if
and only if $X = \Spec A$ is normal and $\Cl X = 0$.
\end{proposition}
